\section{Discussion}\label{sec:discussion}

\subsection{RQ1: what the factor effects mean}

\textbf{The energy result is stronger than the accuracy result.} The MoE uses
42.0\,\% less GPU-board energy per session, and the broader recovered scope
(both GPU boards plus AMD CPU-package energy) still shows a 39.0\,\% reduction.
The latter is a measured-component sensitivity, not full-system energy: the trace
contains no DRAM, board, storage, fan, or power-supply losses. Mean test accuracy
is 2.85\,pp higher for the MoE, but its 95\,\% interval includes zero and the HC3
sensitivity is also nonsignificant. The defensible conclusion is therefore a
large energy advantage with no observed mean-accuracy penalty, not established
accuracy superiority or statistical equivalence.

The mechanism is multiplicative. During proposal windows the MoE device draws
72.5 rather than 112.8\,W and mean proposal latency is roughly half as long.
Whole-session GPU~1 energy is 62.2\,\% lower. GPU~0 session energy also differs
by 11.4\,\% because a board-level counter includes idle time while the proposer
runs; it is not a pure training-phase integral. This distinction matters: separate
devices support direct attribution, but labels such as ``proposer energy'' and
``training energy'' are too strong for their whole-session totals.

The fixed 10- and 20-iteration budgets also narrow the causal claim. A better
proposer cannot reduce the planned number of experiments in this harness; it
changes how many useful mutations those joules produce. The study therefore
demonstrates \emph{energy productivity} (energy per retained improvement), not a
feedback mechanism in which poor steering automatically schedules more training.
A success- or patience-based stopping design would be needed to test that latter
mechanism.

Correctly pooling all sessions, including the six with no retained improvement,
gives 156.4\,kJ/keep for dense and 64.0\,kJ/keep for MoE ($-59.1\,\%$). Five dense
sessions and one MoE session made no retained progress (Fisher $p=0.155$); keep/
revert semantics restore the starting recipe in those cases, so they did not end
below baseline. Candidate quality is nevertheless less stable for dense: its
best evaluated candidate remained below the starting validation baseline in
3/12 sessions, against 0/12 for MoE.

\textbf{Loop budget has diminishing energy productivity.} The original analysis
averaged only defined per-session ratios and silently dropped zero-kept sessions.
The corrected ratio of pooled sums is 61.8 vs.\ 152.1\,kJ/keep for budgets 10 and
20, an increase of 146.2\,\% (bootstrap ratio CI $[1.29,4.46]$). Patience changes
the pooled value by $+24.3\,\%$, not the previously reported $-9.4\,\%$, and its
wide interval includes both directions. These corrections leave the budget
conclusion intact and make the patience conclusion explicitly uncertain.

\subsection{RQ2: reading the frontier}

Two of three sample non-dominated cells are MoE cells, and the cheapest is MoE,
greedy, budget 10 at 56.5\,kJ and 0.8209 mean test accuracy. Energy dominance is
clear; the exact accuracy ordering is not. The dense frontier cell exceeds that
cell by only 0.053\,pp at $1.91\times$ the energy, while the high-budget MoE cell
adds 1.02\,pp at $1.98\times$ the energy. Both accuracy margins are below the
1.23\,pp validation keep threshold, but this threshold measures within-session
re-run noise and is not an equivalence margin for independent test-set means.
Consequently the frontier is practitioner-relevant as a sample summary, with
accuracy membership treated as provisional.

\subsection{Reasoning and temperature follow-ups}

Stage~2a provides exploratory evidence of an architecture--reasoning
interaction. Reasoning changes dense accuracy by $+9.18$\,pp and MoE accuracy by
$-0.91$\,pp; the HC3 interaction is $-10.09$\,pp. Reasoning raises proposer-device
energy in both arms, and the MoE/off cell uses 52.1\,kJ against 172.6\,kJ for
dense-with-reasoning. Their observed accuracies differ by only 0.21\,pp, but the
95\,\% interval ($[-3.97,4.38]$\,pp) is far too wide to establish equivalence.
At $n=3$ per cell, ``reasoning substitutes for sparsity'' is a useful hypothesis
for replication, not a settled generalization.

Stages~2b and 2c establish a narrower temperature result: exact duplicate
proposals are more common at $T=0$ on both arms. They do not establish that
temperature has no effect on retention or accuracy. Each level has only three
sessions, and the Phase-1 restart changed temperature and prompt content together.
The $T=0$ sessions under the current prompt show that deterministic sampling is
not sufficient to force zero retention; they cannot rule out a temperature-by-
prompt interaction. This is why the original diagnosis is reported as unsupported
rather than ``refuted'' in a causal sense.

\subsection{The harness is part of the instrument}

The prompt supplied facts the agent could not otherwise observe (tensor layout,
library versions, and the number of epochs a 45\,s budget permits). Each changed
agent behavior without changing executable workload code. The restart episode
shows why such interventions must change one factor at a time and why early rows
of a randomized run table should not be treated as representative.

An unplanned, functionally equivalent dense/off control occurred in Phase~1 and
Stage~2a. The proposer token cap differed (8,192 vs.\ 12,288) but was non-binding;
session GPU energy agrees to 0.20\,\%, while mean test accuracy differs by
6.42\,pp. Energy measurement is therefore much more stable than the agent-side
accuracy outcome at three repetitions. Follow-ups should allocate more repetitions
to fewer accuracy comparisons.

\subsection{Contract compliance is not proposal quality}

The guard rejected none of the dense proposals and 9.6\,\% of the MoE proposals,
yet the MoE produced the better observed recipes. The pilot showed the same
dissociation. Output-format compliance and proposal quality should therefore be
reported as separate axes rather than combined into one agent score.
